\chapter*{Abstract}
\addcontentsline{toc}{chapter}{Abstract}

Mutual information measures dependence between two variables, but a
difference between two sample estimates need not reflect a difference between
their populations. This thesis studies two independent categorical samples
and tests whether their population mutual information values are equal,
without requiring their complete joint distributions or margins to be the
same.

Normal Wald inference divides a bias-corrected difference by a plug-in
standard error and uses a normal reference. The proposed Expanded
Welch--Satterthwaite procedure keeps that statistic but derives effective
degrees of freedom from the sampling variability of the complete mutual information
variance functional. A change in one table cell also changes the
row and column margins and the pointwise information values; the derivation
accounts for all of these effects. The resulting Student reference is a
moment-matching approximation, not an exact finite-sample distribution.

The completed experiment compares the two methods on fixed positive-support
populations from \(2\times2\) to \(8\times8\), with additional rectangular,
unequal-sample, construction and extreme-sparsity checks. Each of 3,111
unique statistical configurations uses 20,000 paired count-table replicates
at a two-sided nominal level of 0.05. False positives, detection, numerical
validity and runtime are reported separately for individual regimes.
The main null grid includes both identical populations and distinct
populations with equal MI; their calibration is also reported separately.

Expanded Welch moves some liberal Wald false-positive rates toward 0.05, but
makes conservative regimes more conservative and can be less often valid in
sparse samples. Wald is closer to nominal in all 24 main null regimes at
sample sizes 500 or 1,000. At common MI 0.02 and sample size 50,000, both
methods are near nominal in the tested populations; very small MI remains
difficult. Expanded Welch has the same computational order as Wald and costs
about 1.75 times as much per scalar call in the measured implementation.

Near independence, the effective degrees of freedom can become very small,
making Expanded Welch strongly conservative. A supplementary diagnostic also
shows that changing the reference distribution alone cannot correct all
finite-sample errors in the statistic and its standard error.

The evidence supports retaining Wald as the simpler general baseline.
Expanded Welch may be informative in selected liberal regimes, provided that
valid-result frequency and the complete null behaviour are also reported.
